\subsection{Calibration of trade elasticities}\label{app:thetadx}

Our \SOE model is based on the idea that goods are differentiated both on the basis of origin and destination.
In this setting, production supplies goods for the domestic and export markets with a technology capturing the idea that the quantities supplied to these submarkets depend on relative prices.
Each traded good then has three associated trade elasticities: i) $\texttt{etadx(i)}$, the elasticity of transformation between goods produced for the the domestic and export markets; ii) $\texttt{esubdm(i)}$, the elasticity of substitution between domestic and imported varieties entering intermediate and final demand, and iii) $\texttt{esubm(i)}$, the elasticity of substitution between imports from different countries.

In the \GMR model corresponding to the standard \GTAP model, $\eta^{DX} = \infty$, so goods produced for the domestic and export markets are perfect substitutes.
In the \GTAPINGAMS framework, $\eta^{DX}$ can take on any nonnegative value.
A finite value is typically specified in the \SOE model.
When doing this, we adjust the domestic-import demand elasticity to maintain \textit{consistency} with trade responses in the standard \GTAP model.
We interpret consistency from a Marshallian perspective.
In a given market, suppose the import demand elasticity is $\epsilon$,\footnote{When value shares are small, the price elasticity of demand for imports is approximately equal to $-\texttt{esubdm(i)}$.} and benchmark rate of protection is $\bar t$.
Scaling units so that the FOB supply price is $1$, the demand price is $1+\bar t$, and the quantity demanded in the benchmark is $1$.

When the import supply elasticity from partner countries, $\eta$, is infinite, the import supply price is fixed ($p=1$).
Therefore, if the tariff was removed and demand is isoelastic, the new equilibrium quantity and price (in an isoelastic model) are given by:
\begin{equation}
\label{newqty}
q^* = \left(1 + \bar t \right)^{|\epsilon|}\,.
\end{equation}
With a fixed import supply price, the quantity demanded at $p=1$ is the new equilibrium.
This outcome is illustrated in Figure \ref{etainfty}, in which benchmark import demand is at price-quantity point $a$, benchmark import supply is at point $b$, and there is a benchmark trade barrier equal to 20\%.
The free-trade equilibrium is at point $c$, reflecting benchmark import demand elasticity $\epsilon = 5$.

\begin{figure}[p!]

\begin{subfigure}[b]{0.9\textwidth}
\centering
\includegraphics[height=0.4\textheight]{jgea_new/figures/figure9a.tikz}
\caption{ $\epsilon=5$, $\eta=\infty$.}
\label{etainfty}
\end{subfigure}

\vspace{0.2in}

\begin{subfigure}[b]{0.9\textwidth}
\centering
\includegraphics[height=0.4\textheight]{jgea_new/figures/figure9b.tikz}
\caption{$\epsilon=16.7$, $\eta=7.1$.}
\label{theta30}
\end{subfigure}

\caption{Consistent Calibration of $\epsilon$ and $\eta$.}

\end{figure}

Now instead of a perfectly elastic import supply we consider $\eta < \infty$.
In this case, removing tariffs implies a new equilibrium supply price $p^*$, and the associated equilibrium supply and demand $q^*$ from (\ref{newqty}) is given by:
$$q^* = \left( \frac{1+\bar t}{p^*} \right)^{|\epsilon|} =
\left( p^* \right)^\eta\,.$$ The aggregate trade response under $\eta < \infty$ equals the response with $\eta = \infty$ when the demand elasticity ($\epsilon$) is adjusted upward:
$$p^* = (1 + \bar t)^{|\epsilon|/\eta}\,,$$
the reparametrized elasticity of demand is given by:
$$|\epsilon|  = \frac{|\epsilon|}{1 - |\epsilon|/\eta}\,.$$

Let $\theta = 1 - |\epsilon|/\eta$.
Consistency of trade quantity responses in the revised model implies that $\theta \geq 0$.
When $\eta = \infty$, then we have the original model with a horizontal import supply curve.
Setting $\theta=0.3$, the new Marshallian model appears in Figure \ref{theta30}.
As in the previous model, $a$ and $b$ are benchmark demand and supply with a 20\% tariff barrier.
While free trade in the original model is at point $c$, free trade in the revised model is at point $d$.
Supply is less elastic, so social surplus increases are captured by the supply side of the market, and the equilibrium price following tariff removal increases by over 10\% from the benchmark level.

Note that, with this calibration, \texttt{thetadx}$=\theta$ applies in both directions of trade and to all closures: the elasticity of transformation between domestic and export supply is $\eta^{DX}_i=\sigma^D_i/(1-\theta)$ and the elasticity of substitution between domestic and imported goods is raised to $\sigma^{DM}_i=\sigma^D_i/\theta$, where $\sigma^D_i$ is the \GTAP value \texttt{esubd(i)}.
For $\theta=1$ the standard \GTAP specification is obtained ($\eta^{DX}_i=\infty$, $\sigma^{DM}_i=\sigma^D_i$); this is the \texttt{gtap} model version of the present paper.
For $\theta=0.5$, the default of \citet{LanzRutherford2016} and the \texttt{gtapingams} model version of this paper, $\eta^{DX}_i=\sigma^{DM}_i=2\sigma^D_i$.

In the \GAMS code, this adjustment takes place near the end of \texttt{gtap12data.gms}; the default \texttt{thetadx}$=0.5$ is set in \texttt{mge.gms} and may be overridden on the command line.
The run-time environment variable \texttt{thetadx} defines how the adjustment of export supply and import demand is to be calibrated to match the \GMR model.
The associated assignments are reported here.

\begin{Verbatim}
parameter       etadx(g)        Export elasticity of transformation D vs E,
                esubdm(i)       Elasticity of substitution D vs M,
                vxm(g,r)        Export market supply;

$if not set thetadx $set thetadx 0.5

etadx(i)$(1-%thetadx%) = esubd(i)/(1-%thetadx%);
etadx(i)$(not (1-%thetadx%)) = +inf;
esubdm(i) = esubd(i)/%thetadx%;
vxm(i,s) = (vst(i,s) + sum(r,vxmd(i,s,r)))$(1/etadx(i));
\end{Verbatim}

%\clearpage